\documentclass[11pt,fontset=fandol]{ctexart}
\usepackage{reference-answer}
\begin{document}
\papertitle{数学参考答案·排版验收样例}{高中数学\quad 三题版式示例\quad 步骤分值为教学建议。}
\section*{一、基础题}
\question{1}{5}
解方程 \(\sqrt{x+2}=x\)。
\insight{先由根式的非负性限定候选解，再平方求解并代回原式。}
由 \(\sqrt{x+2}=x\) 知 \(x\ge0\)。两边平方，得
\[
 x+2=x^2\quad\Longleftrightarrow\quad(x-2)(x+1)=0.
\]
候选值为 \(x=2\) 或 \(x=-1\)。后者不满足 \(x\ge0\)；前者代回有
\(\sqrt{2+2}=2\)，故原方程的解为 \(x=2\)。
\finalanswer{\(x=2\)。本题按填空题结果给分，解析不另设步骤分。}
\difficulty{基础}{关键是平方前的符号限制与增根检验，运算量小。}

\section*{二、解答题}
\question{2}{15}
某人独立投篮，每次命中概率为 \(p\in(0,1)\)，首次命中或达到 \(N\ge2\) 次时停止，\(X\) 为投篮次数。
\insight{前几次停止需要“先失后中”，最后一次停止只要求此前全失；尾事件更适合直接计算。}
\parttitle{1}{当 \(N=4,p=\frac13\) 时求分布列（建议6分）}
对于 \(j=1,2,3\)，前 \(j-1\) 次未中、第 \(j\) 次命中，故
\[
 P(X=j)=\left(\frac23\right)^{j-1}\frac13.
\]
\checkpoint{2}
而 \(X=4\) 只要求前三次未中，第四次结果不限，故
\(P(X=4)=(\frac23)^3=\frac8{27}\)。\checkpoint{4}
\begin{center}
\begin{tabular}{c@{\qquad}cccc}\toprule
\(X\)&1&2&3&4\\\midrule
\(P\)&\(\frac13\)&\(\frac29\)&\(\frac4{27}\)&\(\frac8{27}\)\\\bottomrule
\end{tabular}
\end{center}
\checkpoint{6}
\parttitle{2}{尾概率与条件概率（建议9分）}
若 \(0\le k\le N-1\)，则 \(X>k\) 等价于前 \(k\) 次均未命中。由独立性，
\[P(X>k)=(1-p)^k.\]\checkpoint{9}
当非负整数 \(k,m\) 满足 \(k+m\le N-1\) 时，
\(\{X>k+m\}\subseteq\{X>k\}\)，且 \(P(X>k)>0\)。因而
\begin{align*}
 P(X>k+m\mid X>k)
 &=\frac{P(X>k+m)}{P(X>k)}\\
 &=\frac{(1-p)^{k+m}}{(1-p)^k}
 =(1-p)^m=P(X>m).
\end{align*}
\checkpoint{15}
\begin{scoring}
建立前三个取值的概率表达式 &2&2\\
正确处理末次投篮的两种结果 &2&4\\
写出完整、数值正确的分布列 &2&6\\
识别尾事件（1分），由独立性计算概率（2分） &3&9\\
事件包含与分母非零（1分），条件概率公式（2分），代入完成证明（3分） &6&15\\
\end{scoring}
\difficulty{中等}{需要区分停止时刻的末项，并把条件事件转化为可计算的尾事件。}

\question{3}{6}
设 \(f(x)=\ln\frac{x}{2-x}-2x+b(x-1)^3\)，\(0<x<2\)。求使 \(f\) 单调递增的 \(b\) 的范围。
\insight{以 \(x=1\) 为中心换元，把导数符号化为含参数的一元不等式。}
令 \(u=x-1\)，则 \(-1<u<1\)，且
\begin{align*}
f'(x)&=\frac2{x(2-x)}-2+3b(x-1)^2\\
&=u^2\left(\frac2{1-u^2}+3b\right).
\end{align*}
\checkpoint{2}
若 \(b<-\frac23\)，则括号内的式子在 \(u=0\) 附近为负，故 \(f'\) 在某个非空区间为负，不能递增。
\checkpoint{4}
若 \(b\ge-\frac23\)，则当 \(u\ne0\) 时
\[
\frac2{1-u^2}+3b\ge\frac2{1-u^2}-2=\frac{2u^2}{1-u^2}>0.
\]
故 \(f'(x)>0\) 对 \(x\ne1\) 成立，而 \(f'(1)=0\)，因此 \(f\) 在 \((0,2)\) 上严格递增。
\finalanswer{\(b\in[-\frac23,+\infty)\)。}\checkpoint{6}
\begin{scoring}
正确求导、换元并因式分解 &2&2\\
利用零点附近符号证明必要性 &2&4\\
证明充分性并给出参数范围 &2&6\\
\end{scoring}
\difficulty{较难}{需区分必要性与充分性，并处理孤立驻点不破坏严格单调性。}
\end{document}
